\documentclass[10pt,a4paper]{amsart}
\usepackage[margin=19mm]{geometry}
\usepackage{amsmath,amssymb,mathtools}
\usepackage[scr=rsfs,cal=euler]{mathalfa}
\usepackage{xfrac}
\usepackage[unicode,hidelinks,bookmarks=false]{hyperref}
\newcommand{\ie}{i.e.\ }
\newcommand{\cf}{cf.\ }
\newcommand{\resp}{respectively\ }
\newcommand{\Subsection}{\subsection}
\providecommand{\nobreakdash}{\nobreak\hbox{-}}
\setlength{\emergencystretch}{2em}
\multlinegap0pt
\usepackage{bm}
\usepackage{bbm}
\usepackage{dsfont}
\usepackage{xspace}
\usepackage{cancel}
\usepackage{mathtools}
\usepackage{amsthm}
\makeatletter
\@ifundefined{theorem}{\newtheorem{theorem}{Theorem}}{}
\@ifundefined{lemma}{\newtheorem{lemma}[theorem]{Lemma}}{}
\@ifundefined{proposition}{\newtheorem{proposition}[theorem]{Proposition}}{}
\makeatother
\newcommand{\cA}{\mathcal{A}}
\newcommand{\cB}{\mathcal{B}}
\newcommand{\cC}{\mathcal{C}}
\title[Binomiality of colored Gaussian models]{Binomiality of colored Gaussian models}
\author{Benjamin Biaggi, Jan Draisma, Magdal\'ena Mi\v{s}inov\'a}
\date{Source: arXiv:2507.06437v1 (CC BY 4.0)}
\begin{document}
\maketitle

\noindent\emph{Source and licence.} Binomiality of colored Gaussian models. Version arXiv:2507.06437v1 (CC BY 4.0), distributed under the Creative Commons Attribution 4.0 International licence, as recorded in the arXiv licence metadata for this exact version and shown on the abstract page https://arxiv.org/abs/2507.06437v1. Full rights notice: licenses/arxiv-2507.06437-CC-BY-4.0.txt.\par\smallskip

\noindent\textit{Editorial scope.} Counted unit: the Theorem whose source label is \texttt{thm:sets-of-paths} in the article (source0.tex lines 1519--1528, source label \texttt{thm:sets-of-paths}), delivered together with the article's own complete original proof of it (source0.tex lines 1530--1603, one \texttt{proof} environment of 74 lines). No mathematics is written, repaired or re-derived: statement and proof are the article's own text line for line. Required context retained uncounted: prop:removing-cliques (lines 982--1002); lemma:moving-paths (lines 1237--1245); lemma:symmetry (lines 1297--1300). Every definition, display and label of those lines is retained unchanged. Omitted from this package and not redistributed: 1--981, 1003--1236, 1246--1296, 1301--1518, 1529--1529, 1604--1633. Nothing inside a retained excerpt is omitted or altered: every retained body is delivered exactly as the article's lines are written, so no omission receipt of any kind accompanies this package. Citations: the retained excerpts carry no citation command at all, so no bibliography entry of the article is transcribed and no reference is redistributed. Reference adaptation: none was needed and none was made; every reference inside the delivered unit resolves inside it, so no unresolved reference or citation is produced. Printed slips kept literal: no printed word, symbol, hypothesis or number of the counted statement or of its proof was altered. Layout and class: the article's own class is \texttt{amsart} with packages including graphicx, babel, amssymb, amsmath, amsthm, geometry, hyperref, bbm, booktabs, tikz-cd, enumitem, amscd, mathtools; the wrapper is the lane's pinned \texttt{amsart} editorial wrapper at 10pt on A4, which restates the theorem-like environments the retained text needs and adds the article's own macro definitions that the retained text needs (\texttt{\textbackslash cA}, \texttt{\textbackslash cB}, \texttt{\textbackslash cC}), with extra wrapper packages (bm, bbm, dsfont, xspace, cancel, mathtools). The delivered numbering is the unit's own. Assets: no retained span contains a figure environment, a graphics-inclusion command, a PDF-page inclusion, a table or a TikZ picture; no image file is copied into this package, the package declares no asset dependency and no image file is redistributed with it. Licence: version arXiv:2507.06437v1 of this article is distributed under the Creative Commons Attribution 4.0 International licence, as recorded in the arXiv licence metadata for this exact version and shown on the abstract page https://arxiv.org/abs/2507.06437v1. A full rights notice is delivered at licenses/arxiv-2507.06437-CC-BY-4.0.txt. Characters: every retained excerpt is pure ASCII, so the delivered engine, XeLaTeX with its default Latin Modern text font, reports no missing character.
\begin{proposition}\label{prop:removing-cliques}
    Given a connected, triangle-regular block graph $G$ with at least two maximal
    cliques, we can remove a set $\mathcal C$ of cliques from $G$ so that the
    following conditions hold. Here, by removing we mean deleting all
    the vertices contained {\em only} in cliques in $\mathcal{C}$. 
    \begin{enumerate}
        \item We obtain a connected triangle-regular block graph.
        \item All the cliques in $\mathcal{C}$ have the same multiset
	of colored vertices and colored edges.
        \item If $v\in C\in \mathcal{C}$ is a vertex, and $v' \in V(G)$ is a
	vertex of the same color: $\lambda(v')=\lambda(v)$, then $v'\in V(C')$ for some $C'\in \mathcal{C}$, similarly for edges.
        \item For each $C\in \mathcal{C}$ there is exactly one vertex
	$v_C\in V(C)$ contained in at least two maximal cliques. The
	color $\lambda(v_C)$  is independent of the choice of $C$, and
	the only vertices of $G$ with this color are the vertices
	$v_C$ with $C$ ranging through $\cC$.
        \item For any vertex $v\in V(C)\setminus \lbrace v_C\rbrace$,
	$\lambda(\{v,v_C\})$ depends only on $\lambda(v)$.
        \item Let $C\in \mathcal C$ and let $v_{C'}$ be a vertex with $\lambda(v_{C'}) = \lambda(v_C)$. Then the number of cliques in $\mathcal C$ containing $v_{C'}$ is the same as for $v_C$.
    \end{enumerate}
\end{proposition}

\begin{lemma}\label{lemma:moving-paths}
    In a connected, triangle-regular block graph $G$, let $P$ and $Q$ be
    any two isomorphic shortest paths, and let $P'$ be a shortest path
     containing $P$. Then there exists a quasi-automorphism
    $\alpha$ of $G$ such that $\alpha(P) = Q$ (and indeed, $\alpha$ can
    be chosen to agree with any of the isomorphisms between $P$ and $Q$,
    of which there may be two) and $\lambda(\alpha(e)) = \lambda(e)$
    for all edges $e$ in $P'$.
\end{lemma}

\begin{lemma}\label{lemma:symmetry}
    For any vertices $u,v\in G$ of the same color, $u\leftrightarrow
    v$ is palindromic.
\end{lemma}

\begin{theorem}\label{thm:sets-of-paths}
    Let $\mathcal A$ and $\mathcal B$ be two multisets of shortest paths in a triangle-regular block graph $G$ such that the multiset of colors of edges and the multiset of colors of vertices in $\mathcal{A}$ are the same as those in $\mathcal B$. We are allowed to modify those two sets by the following two moves:
    \begin{enumerate}
        \item Replace a path in $\cA$ by an isomorphic shortest path.
        \item Replace a pair $(P,Q)$ of shortest paths in $\cA$ that
	meet at some vertex $x$ by the swap $((P_1,Q_2),(Q_1,P_2))$ of
	$(P,Q)$ at $x$.
    \end{enumerate}
    Then there exists a finite sequence of moves which achieves $\mathcal A=\mathcal B$.
\end{theorem}

\begin{proof}
We proceed by a double induction: on the cardinality of $\cA$, and then
on the sum of the lengths of the paths in $\cA$.  If all paths in $\cA$
have length zero, then the multiset of edge colors of paths in $\cA$
is empty and hence the same holds for $\cB$. In this case, the theorem
clearly holds. 

Otherwise, let $k$ be the maximum depth among all edges appearing in
paths of $\cA$. Such an edge either has two vertices of depth $k$ or
else one vertex of depth $k$ and one vertex of depth $k-1$.

If the former case applies to some $e$, then any path $P$ in
$\cA$ containing a copy of $e$ consists of $e$ and its vertices
only. By the assumptions on $\cA$ and $\cB$, $\cB$ has a path
$P'$ that contains an edge $e'$ with $\lambda(e')=\lambda(e)$. By
Proposition~\ref{prop:removing-cliques}, $P'$ equals $(u',f',v')$ and
is therefore isomorphic to $P$. The theorem now reduces to the smaller
instance with $\cA \setminus \{P\}$ and $\cB \setminus \{P'\}$, and we
are done by induction.

Suppose that there strictly more vertices of depth $k$ in paths in $\cA$
than there are edges of this depth. Then some path $P$ in $\cA$ consists
of a single depth-$k$ vertex $u$. The assumptions on $\cB$ and $\cA$
imply that among the vertices in paths in $\cB$ there are more of color
$\lambda(u)$ than there are edges in paths in $\cB$ with an endpoint of
color $\lambda(u)$. Hence $\cB$, too, has a path $P'$ consisting of a
single vertex $u'$ with $\lambda(u)=\lambda(u')$. Removing $P$ from $\cA$
and the isomorphic path $P'$ from $\cB$, we are done by induction.

So we may assume that all edges of depth $k$ in paths in $\cA$ connect
a vertex of depth $k$ to a vertex of depth $k-1$, and that all paths in
$\cA$ with a vertex $u$ of depth $k$ have $u$ as endpoint and contain
at least one edge, connecting $u$ to a vertex of depth $k-1$. The same
statements then apply to $\cB$. 

Pick an edge $e$ of depth $k$ in a path $P$ in $\cA$, and let $u,v$
be its vertices of depth $k,k-1$, respectively. Then $\cB$ contains
a path $P'$ with an edge $e'$ with $\lambda(e')=\lambda(e)$, and it
follows that $\kappa(e')=k$, the vertex $u'$ of $e'$ of depth $k$ is
the endpoint of $P'$, and the other vertex $v'$ of $e'$ has depth $k-1$.

Now remove $e$ and $u$ from $P$ to obtain $\Tilde{P}$ and $e$ and $u'$
from $P'$ to obtain $\Tilde{P}'$, and let $\Tilde{\cA}$ and $\Tilde{\cB}$
be the shortest path collections obtained by replacing $P$ in $\cA$
by $\Tilde{P}$ and $P'$ in $\cB$ by $\Tilde{P}'$. As the sum of the
lengths of paths in $\cA$ has shrunk by $1$, the induction
hypothesis says that we can make $\Tilde{\cA}$ equal to $\Tilde{\cB}$ by a sequence
of moves.

We carry out the same moves for $\cA$ in the following sense. In moves
of type (1), whenever $\Tilde{P}$ is replaced by an isomorphic path
$\Tilde{Q}$, by Lemma~\ref{lemma:moving-paths} there is a shortest path
$Q$ in $G$ isomorphic to $P$ containing $\Tilde{Q}$, and we replace $P$
in $\cA$ by $Q$. In the process we update $e$ and $u$ and $v$
accordingly, so that $P$ still starts with $(u,e,v)$. And in a move of type (2)
that involves $\Tilde{P}$ and another path $Q$, say with $v$ an endpoint of $\Tilde{P}_i$, we update $\Tilde{P}$ to the new path containing $\Tilde{P}_i$,
so that $\tilde{P}$ again has $v$ as one endpoint, and in $\cA$ we 
replace $(P,Q)$ by the swap of $(P,Q)$ at $x$; then the new $P$ is the
new $\tilde{P}$ with $e$ and $u$ attached. 

After these moves, $\cA$ and $\cB$ become equal as multisets when
we remove $u,e$ from $P$ and $u',e'$ from $P'$; we continue to write
$\Tilde{P}$ and $\Tilde{P}'$, respectively, for these truncated paths.
If $\Tilde{P}=\Tilde{P}'$, then $P$ and $P'$ are isomorphic (this needs
Lemma~\ref{lemma:symmetry} when both endpoints of $\Tilde{P}$ have the
color $\lambda(v)$). In that case, we can remove $P$ from $\cA$ and $P'$
from $\cB$, and we are done by induction.  Otherwise, let $Q \in \cA$
be an element equal to $\Tilde{P}'$. In particular, $Q$ has $v'$ as an
endpoint.  Replace $P \in \cA$ by an isomorphic path that contains
$u',e',v'$. Then replace $(P,Q)$ in $A$ by its swap at $v'$; this has
the effect that $u'$ and $e'$ are removed from $P$ and attached to $Q$,
so that $Q$ becomes equal to $P'$. After removing $Q$ from $\cA$ and $P'$
from $\cB$, we are done by induction.
\end{proof}

\end{document}
